% Debug report: the problems this build actually hit and how each was resolved.
% Built with `make report-debug`. Equal rank deliverable to report/main.pdf.
\documentclass[11pt,a4paper]{article}

\usepackage[T1]{fontenc}
\usepackage[utf8]{inputenc}
\usepackage{lmodern}
\usepackage[margin=2.6cm]{geometry}
\usepackage{booktabs}
\usepackage{microtype}
\usepackage{amsmath}
\usepackage{xcolor}
\usepackage{fancyvrb}
\usepackage[hidelinks]{hyperref}

\definecolor{rulegrey}{HTML}{C3C2B7}
\setlength{\parskip}{0.45em}
\setlength{\parindent}{0pt}

\newcommand{\entry}[2]{%
  \vspace{0.6em}%
  {\color{rulegrey}\hrule height 0.6pt}%
  \vspace{0.5em}%
  \subsection*{#1}%
  \textit{\small #2}\par\vspace{0.3em}%
}
\newcommand{\field}[1]{\vspace{0.35em}\textbf{#1}\quad}

\title{\textbf{ML Experiment Triage: Debug Report}\\[0.4em]
\large What went wrong, what caused it, and how each was resolved}
\author{Olajide Badejo}
\date{5 August 2026}

\begin{document}
\maketitle

\section*{Why this document exists}

The main report says what the tool does and how well it works. This one says what
it took to get there, which is the part that is normally thrown away and is
usually the part with the most in it.

Nine entries follow. Three are defects in my code, two are cases where my test
was wrong and the code was right, one is a flaw in my own measurement method, one
is a toolchain misconfiguration that silently disabled the driver I had
specified, one is a bug caught by reading before it could produce a wrong answer,
and one is an environment decision. They are in the order they were hit.

The one worth reading if you read only one is the third: a statistical test that
was running at better than twice its nominal error rate while every mechanics
test around it passed.

% ---------------------------------------------------------------------------
\section{Environment}

\entry{No Python 3.13 on the machine}{Symptom: \texttt{py -0p} listed only 3.14,
where the build targets 3.13.}

\field{Root cause} Nothing broken. The machine simply had a newer interpreter
than the project targets.

\field{Options} Build against 3.14 and see what breaks; install 3.13; fall back
to 3.12.

\field{Fix and why} Installed 3.13.14 with \texttt{py install 3.13}. Building
against 3.14 would have meant betting on the wheel situation for tensorboard and
protobuf, the two dependencies here most likely to lag a new interpreter, and any
failure would then have been ambiguous between my code and a missing wheel.

\field{Verification} A TensorBoard write and read round trip was run before any
parser code existed, so that a later parser failure could not be blamed on the
toolchain.

% ---------------------------------------------------------------------------
\section{Tooling}

\entry{The dash guard flagged a file that is not in the repository}{Symptom: the
first run failed on an em dash in the build specification sitting in the working
directory.}

\field{Root cause} The guard walked the filesystem. The file it flagged is
ignored and is an input to the work rather than part of the project, so
filesystem scope was answering the wrong question. The same bug would later have
flagged anything inside the virtual environment or a build directory.

\field{Options} Add path exclusions and keep adding them forever; or scope the
guard to what git would ship.

\field{Fix and why} The guard now asks
\verb|git ls-files --cached --others --exclude-standard|, which is exactly the
set of files that are or will become part of the repository. Exclusion lists rot;
asking git does not.

\field{Verification} Passes on a clean tree, still fails on a planted dash in a
tracked file, and runs twice in \texttt{make all}, once before the PDFs exist and
once after, so the compiled PDFs are checked too.

\entry{texify was never running}{Symptom: \texttt{make report} produced a correct
PDF, but the log showed \texttt{-{}-batch-mode: unknown option} followed by a
fallback to a manual pdflatex sequence.}

\field{Root cause} MiKTeX \texttt{texify} spells the flag \texttt{-{}-batch}.
\texttt{-{}-batch-mode} is pdflatex's spelling, and I had written the pdflatex
one.

Two further lessons hide in this entry, both about writing the entry itself.
Typing those flag names as a literal pair of hyphens in LaTeX prose typesets an
en dash, and the project's dash guard caught it in the source \emph{and} in the
compiled PDF, which is the guard doing exactly the job it was built for on the
document describing the guard. The obvious repair, wrapping them in
\texttt{\textbackslash verb}, then failed to compile at all, because
\texttt{\textbackslash verb} is illegal inside a macro argument and these sit
inside one. The idiom that works in both places is \texttt{-\{\}-}, which breaks
the ligature without needing verbatim mode.

\field{Why it mattered} The fallback path existed for continuous integration,
where texify does not exist, and it worked, so the build was green and the
specified driver was silently never used. A fallback that quietly becomes the
primary path is worse than no fallback, because it hides the fact that the
configuration is wrong.

\field{Fix and why} Corrected the flag. The fallback stays, because the Ubuntu
runner genuinely needs it, but it now announces itself when it triggers.

\field{Verification} The build log names the driver it used, and now shows
texify.

% ---------------------------------------------------------------------------
\section{Statistics}

\entry{The window block mode was running at 12 percent type I error}{Symptom:
with a nominal alpha of 0.05 and a true effect of exactly zero, the single run
mode rejected the null on 12.0 percent of cases, against a gate of [2, 8]
percent.}

The seed replicated mode passed on its first run at 4.6 percent, so the defect
was specific to the block machinery. Every mechanics test around it passed. This
was found only because the calibration suite exists.

\field{Investigation} A diagnostic printed block count, block length and
autocorrelation time per case, then swept type I error against window length,
autocorrelation level, block length multiplier, and whether the statistic was
smoothed. Two causes fell out of that grid.

\field{Root cause, part one} The autocorrelation time was estimated on the two
windows \emph{concatenated}. Joining two independent runs end to end puts a step
change at the join, and the estimator reads a step change as long range
dependence. Measured, it inflated the estimate by roughly a factor of two. It was
also masked: the inflated value was then clipped by a target block count, so the
net effect was neither reliably conservative nor reliably wrong, which is the
worst kind of bug to have in a statistic.

\field{Root cause, part two} The block length was one autocorrelation time. At
one $\tau$ the block means are still visibly correlated with their neighbours.
The observed split is contiguous, one run's blocks against the other's, but a
permuted split scatters neighbouring blocks across both groups, breaking that
correlation and producing a null distribution narrower than the truth. A null
that is too narrow is precisely how a test becomes anticonservative.

Smoothing made this worse in a way I had not anticipated. The nine point moving
average the statistic uses raises the measured autocorrelation time of the window
from about 8 to about 14, so the blocks needed to be longer than the raw series
would have suggested.

\field{Options} (a) Keep one $\tau$ and accept the error rate. Rejected: the
whole claim of the project is that these p values mean something. (b) Raise the
multiplier and clip the block length when the window is too short. Rejected after
measuring it: clipping reintroduces exactly the failure, and the measured rate at
short windows stayed above 10 percent. (c) Raise the multiplier and refuse when
the window cannot support it. (d) Abandon the single run mode. Rejected:
sometimes one run is all there is, and a labelled weak answer beats no answer.

\field{Fix and why} Option (c). $\tau$ is now estimated per window with the larger
taken, the block is three times $\tau$, and below eight blocks the comparison
raises an error naming the concrete remedies rather than returning a number.
Three was chosen by measurement: across window lengths from 1200 to 8000 steps
and autocorrelation coefficients from 0 to 0.9, one $\tau$ gave 9 to 14 percent,
two gave 7 to 10, and three landed on nominal wherever the precondition was
comfortably satisfied.

Refusing rather than clipping is the part I would defend hardest. A quietly
miscalibrated p value is worse than no p value, because it spends credibility
that nothing later can recover.

\field{Verification} Type I error 5.81 percent over 1928 null cases, 72 refused
as too short. Power 100 percent on a large effect, confirming the fix did not
simply make the test deaf. A separate test pins the refusal.

\field{Commit} Phase 2.

\entry{A selection effect in my own measurement}{Symptom: after the fix above,
the measured type I error still looked high in some cells: 20 percent at 3000
steps with an autocorrelation coefficient of 0.9, against 5.6 percent at 8000
steps.}

\field{Root cause} Not the estimator. The refusal rule decides using $\tau$
estimated from the same data it then tests. In the cells where most cases were
refused, the survivors were exactly those whose $\tau$ happened to be
underestimated, which are exactly the cases that get blocks that are too short.
Conditioning on survival selected for the failure. Every cell with a refusal rate
under 10 percent measured between 4.4 and 5.8 percent, and the error rate rose
monotonically with the refusal rate, which is the signature of a selection effect
rather than a broken statistic.

\field{Options} Refuse using something independent of the data, which for an
unknown autocorrelation time is not available; report the conditional rate and
say so; or run the suite inside the region where the precondition is comfortably
satisfied.

\field{Fix and why} The last, with the second written down. The suite runs at
8000 steps, where under 4 percent of cases are refused, and it asserts that the
refusal rate stays below 10 percent so the measurement cannot silently drift into
the biased regime. The honest statement, which is in the methodology, is that the
measured rate is conditional on the precondition being met.

\field{Verification} The suite asserts both the error rate and the refusal rate.

% ---------------------------------------------------------------------------
\section{Reproducibility}

\entry{Every HTML report differed from the last}{Symptom: a test comparing two
renderings from the same database at the same seed failed immediately on being
written.}

\field{Root cause} Plotly assigns each figure a freshly generated identifier when
rendering to HTML. Nothing about the data changed; the container changed.

\field{Why it mattered more than it looked} The design says reports are pure
functions of the database plus a recorded seed, and the README says so. A file
that differs on every run makes that claim unverifiable: nobody can diff two
reports to confirm nothing moved, which is the practical way anyone would ever
check it.

\field{Options} Strip the generated identifiers after rendering, which is fragile
against the library's own internal references to them; or set them explicitly.

\field{Fix and why} Set explicitly from the metric name, which is already unique
per figure and is more readable in the output.

\field{Verification} Two renderings now agree line for line, ignoring only the
generation timestamp, which is the one thing that legitimately differs.

\entry{String hashing would have broken reproducibility silently}{Symptom: none.
Caught by reading, before it could produce a wrong result.}

\field{Root cause} The synthetic sweep seeded each run's generator with
\verb|hash(condition.name)|. Python randomises string hashing per process unless
\verb|PYTHONHASHSEED| is set, so the sweep would have differed on every
invocation while looking entirely deterministic in the code.

\field{Why it would have been nasty} Every test in the repository would still
have passed, because they all regenerate the sweep within a single process. The
failure would have surfaced only as the committed verdict record never matching,
and the cause would have been very hard to see from that symptom.

\field{Fix and why} \verb|zlib.crc32| of the name, which is stable across
processes and versions. Keying on the name rather than an index also means adding
a condition later cannot shift the numbers of the conditions before it.

\field{Verification} \texttt{make verify-demo} rebuilds the sweep from its seeds
and checks all twelve verdicts against the committed record to four decimal
places.

% ---------------------------------------------------------------------------
\section{Where the test was wrong and the code was right}

Recorded because being wrong in this direction is worth noticing too. In both
cases the repair was to assert the real property, not to loosen the assertion.

\entry{The smoothing bleed}{Symptom: a test built a step function, 900 zeros then
100 fives, and asserted the window statistic was 5.0 within 0.05. It measured
4.9444.}

\field{Root cause} The nine point moving average reaches back across the step, so
the four points at the start of the window each borrow some of the value before
it. The deficit is $4 \times (4+3+2+1) / 9$ spread over 100 points, which is
exactly 0.0556. The code was right; my tolerance was wrong.

\field{Fix and why} The test now asserts the derived value to six decimal places,
and a second test covers the clean case where nothing bleeds. That is a stronger
check than the loose one it replaced.

\entry{The self containment check}{Symptom: a test asserting that no external URL
appears in the HTML report failed.}

\field{Root cause} The inlined Plotly bundle carries a default topojson host and
several map tile attributions as string literals, in code paths a line chart
never touches. None of it is fetched.

\field{Fix and why} The test now strips the inlined scripts and checks the
remaining markup for tags that actually cause a network request. A test that
fails on correct code is not a strict test, it is a broken one.

\entry{I was wrong about the learning rate}{Symptom: a test asserting that
learning rate is the strongest correlate of the validation loss failed. Batch
size correlated more strongly.}

\field{Root cause} Not a bug. The sweep gives learning rate an optimum in the
middle of its range: 0.0003 is worse than 0.001, 0.003 is the best of the four,
and 0.01 is by a distance the worst. That is a U shape, and Spearman rank
correlation measures monotone association, so the correlation of a U shape is
near zero however large the effect. Batch size was swept monotonically over the
same runs, so it correlates more strongly while driving a smaller effect.

\field{Options} (a) Change the demo sweep so learning rate is monotone and the
test passes. Rejected, and it is worth saying why: it would have made the demo
flatter the method by hiding the exact case where the method is blind, and a
learning rate with an optimum is the realistic shape, not the awkward one.
(b) Add a nonlinear dependence measure. Deferred as scope, recorded as future
work. (c) Keep the behaviour, document the limitation, pin it with a test.

\field{Fix and why} Option (c). The limitation is now in the module docstring,
the methodology, the README, the main report and the caption under the
sensitivity table in the HTML. The test was renamed for the blindness it checks,
so that a future reader who sees a weak learning rate correlation does not
``fix'' it. A second assertion confirms the effect the correlation cannot see is
genuinely there and correctly sized: the tool recovers a gap of 0.193 between the
best and worst conditions against a designed 0.190.

\field{Verification} Two tests, one at unit level on an isolated U shape and one
at integration level on the real sweep.

% ---------------------------------------------------------------------------
\section{What I take from this}

Three things.

\textbf{The calibration suite paid for the whole project in one run.} It found a
statistical test running at better than twice its nominal error rate, in code
that looked correct and passed every mechanics test written for it. No amount of
code review would have found that number. Roughly a minute of measurement did.

\textbf{Refusing is a feature.} Two of the fixes above ended in the tool
declining to answer: below eight blocks it raises rather than returning a p
value, and below the attainable alpha it reports inconclusive rather than no
change. Both were more work than the alternative and both are the reason the
error rates hold.

\textbf{A failing test is a question, not an instruction.} Three of the nine
entries here are cases where the test was wrong. In each, the temptation was to
adjust the assertion until it passed, and in each the right move was to work out
what the code was actually doing and then assert that precisely. One of them
turned a bug report into a documented limitation of the method, which is more
useful to a reader than the clean result I had expected would have been.

\end{document}
